\documentclass[10pt,a4paper]{amsart}
\usepackage[margin=19mm]{geometry}
\usepackage{amsmath,amssymb,mathtools}
\usepackage[scr=rsfs,cal=euler]{mathalfa}
\usepackage{xfrac}
\usepackage[unicode,hidelinks,bookmarks=false]{hyperref}
\newcommand{\ie}{i.e.\ }
\newcommand{\cf}{cf.\ }
\newcommand{\resp}{respectively\ }
\newcommand{\Subsection}{\subsection}
\providecommand{\nobreakdash}{\nobreak\hbox{-}}
\setlength{\emergencystretch}{2em}
\multlinegap0pt
\usepackage{bm}
\usepackage{bbm}
\usepackage{dsfont}
\usepackage{xspace}
\usepackage{cancel}
\usepackage{mathtools}
\usepackage[shortlabels]{enumitem}
\usepackage{amsthm}
\makeatletter
\@ifundefined{theorem}{\newtheorem{theorem}{Theorem}}{}
\makeatother
\newcommand\mb{\mathbb}
\newcommand{\F}{\mathcal{F}}
\newcommand{\R}{\mathbb{R}}
\title[Dangerous Liaisons of Convex Learning and Non-Affine Aggrega]{Dangerous Liaisons of Convex Learning and Non-Affine Aggregation}
\author{Thomas Boudou, Batiste Le Bars, Nirupam Gupta, Aur\'elien Bellet}
\date{Source: arXiv:2606.28123v1 (CC BY 4.0)}
\begin{document}
\maketitle

\noindent\emph{Source and licence.} Dangerous Liaisons of Convex Learning and Non-Affine Aggregation. Version arXiv:2606.28123v1 (CC BY 4.0), distributed under the Creative Commons Attribution 4.0 International licence, as recorded in the arXiv licence metadata for this exact version and shown on the abstract page https://arxiv.org/abs/2606.28123v1. Full rights notice: licenses/arxiv-2606.28123-CC-BY-4.0.txt.\par\smallskip

\noindent\textit{Editorial scope.} Counted unit: the Theorem whose source label is \texttt{thm:universal\_convexity} in the article (source0.tex lines 351--359, source label \texttt{thm:universal\_convexity}), delivered together with the article's own complete original proof of it (source0.tex lines 360--501, one \texttt{proof} environment of 142 lines). No mathematics is written, repaired or re-derived: statement and proof are the article's own text line for line. Required context retained uncounted: eq:algo-update (lines 226--229); eq:operator-convexity (lines 247--249). Every definition, display and label of those lines is retained unchanged. Omitted from this package and not redistributed: 1--225, 230--246, 250--350, 502--3891. Nothing inside a retained excerpt is omitted or altered: every retained body is delivered exactly as the article's lines are written, so no omission receipt of any kind accompanies this package. Citations: the retained excerpts carry no citation command at all, so no bibliography entry of the article is transcribed and no reference is redistributed. Reference adaptation: none was needed and none was made; every reference inside the delivered unit resolves inside it, so no unresolved reference or citation is produced. Printed slips kept literal: no printed word, symbol, hypothesis or number of the counted statement or of its proof was altered. Layout and class: the article's own class is \texttt{article} with packages including jmlr2e, lastpage, mathtools, mathrsfs, booktabs, hyperref, comment, enumitem, bibentry, mdframed, tikz, cleveref, tcolorbox, amsmath; the wrapper is the lane's pinned \texttt{amsart} editorial wrapper at 10pt on A4, which restates the theorem-like environments the retained text needs and adds the article's own macro definitions that the retained text needs (\texttt{\textbackslash F}, \texttt{\textbackslash R}, \texttt{\textbackslash mb}), with extra wrapper packages (bm, bbm, dsfont, xspace, cancel, mathtools, enumitem). The delivered numbering is the unit's own. Assets: no retained span contains a figure environment, a graphics-inclusion command, a PDF-page inclusion, a table or a TikZ picture; no image file is copied into this package, the package declares no asset dependency and no image file is redistributed with it. Licence: version arXiv:2606.28123v1 of this article is distributed under the Creative Commons Attribution 4.0 International licence, as recorded in the arXiv licence metadata for this exact version and shown on the abstract page https://arxiv.org/abs/2606.28123v1. A full rights notice is delivered at licenses/arxiv-2606.28123-CC-BY-4.0.txt. Characters: every retained excerpt is pure ASCII, so the delivered engine, XeLaTeX with its default Latin Modern text font, reports no missing character.
\begin{equation}\label{eq:algo-update}
    % G^F_{\gamma}(\theta_t)
    \theta_{t+1} = \theta_t - \gamma F(g_1(\theta_t), \ldots, g_n(\theta_t)), \text{ where } \theta_0 \in \R^d, \gamma > 0.
\end{equation}

\begin{equation}\label{eq:operator-convexity}
    \forall \theta, \omega \in \R^d,\quad \langle \theta - \omega, \F(\theta) - \F(\omega) \rangle \geq 0.
\end{equation}

\begin{theorem}\label{thm:universal_convexity}
    Let $d>1$, and consider the model update rule~\eqref{eq:algo-update}. 
    $\F$ is a monotonic operator for all gradients $\{g_i(\cdot)\}_{i\in[n]}$ induced by convex functions
    if and only if $F$ is a positively affine aggregation, i.e., there exists $\alpha \in \mb{R}_+^n$ and $C \in \mb{R}^d$, such that for any input vectors $(g_1, \ldots, g_n) \in (\mb{R}^d)^n$,
    \[
        F(g_1, \ldots, g_n) = \sum_{k=1}^n \alpha_k g_k + C.
    \]
    If $F$ is also permutation-invariant and idempotent ($F(g, \ldots,g)\! =\! g$), then $F$ is the arithmetic mean.%\looseness=-1
\end{theorem}

\begin{proof}
    If $F$ is a positively affine aggregation, it is immediate that $\F$ preserves monotonicity.
    We now turn to the converse. 
    The proof has three steps.
    First, Step 1 proves a rigid alignment condition: any variation in a single input vector must yield a variation of the aggregation output in the exact same direction.
    We prove this by contradiction, showing that if this property is violated, one can construct valid inputs $\{g_i(\cdot)\}_{i\in[n]}$ such that monotonicity~\eqref{eq:operator-convexity} is violated.
    Then, Steps 2 and 3 deduce the necessary algebraic form of $F$ from this constraint, proceeding first input-wise and then globally.

    \paragraph{Step 1.} We prove that any variation in a single input vector must result in an output shift that is a non-negative scalar multiple of that variation. 
    Fix an index $i \in [n]$ and an arbitrary vector $g_{-i} \vcentcolon= (\ldots, g_{i-1}, g_{i+1}, \ldots) \in (\R^d)^{n-1}$. 
    Consider the partial function
    \[
        f_{g_{-i}}(u) := F(g_1, \ldots, u, \ldots, g_n).
    \]
    Let $u, v \in \mathbb{R}^d$ be arbitrary vectors. 
    We claim that if~\eqref{eq:operator-convexity} is satisfied,
    \[
        f_{g_{-i}}(u) - f_{g_{-i}}(v) \in \operatorname{cone}(u-v) %\vcentcolon= \{\lambda(u-v) ; \lambda \geq 0\}.
    \] 
    If $u = v$, or $f_{g_{-i}}(u) = f_{g_{-i}}(v)$, the claim is trivial. 
    Otherwise, the proof proceeds by contradiction. Suppose that $f_{g_{-i}}(u) - f_{g_{-i}}(v) \notin \operatorname{cone}(u-v)$.
    We define $\delta$ as the bisector of the angle between $(u-v)$ and $-(f_{g_{-i}}(u) - f_{g_{-i}}(v))$,
    \[
        \delta \vcentcolon= \frac{u-v}{\|u-v\|_2} - \frac{f_{g_{-i}}(u) - f_{g_{-i}}(v)}{\|f_{g_{-i}}(u) - f_{g_{-i}}(v)\|_2}.
    \]
    %By the hyperplane separation theorem~\citep[separating a point from a half line,][Theorem 11.4]{rockafellar1970convex}, there exists $\delta \in \mathbb{R}^d$ such that
    As directions differ, the strict Cauchy-Schwarz inequality applies and implies the separation 
    \begin{equation}\label{eq:strict_separation}
        \langle f_{g_{-i}}(u) - f_{g_{-i}}(v), \delta \rangle < 0 \quad \text{and} \quad \langle u-v, \delta \rangle > 0.
    \end{equation}
    Next, we construct convex functions whose gradients at $0$ and $\delta$ match these input configurations.
    \begin{enumerate}[label=$\spadesuit$, ref=($\spadesuit$)]
        \item\label{item:def_R_i} For index $i$, define $\ell_i(\theta) = \frac{1}{2} \theta^\intercal H \theta + v^\intercal \theta$, with $H = \frac{(u-v) (u-v)^\intercal}{\langle u-v, \delta \rangle}$. 
        Since $\langle u-v, \delta \rangle > 0$, $H_i$ is symmetric p.s.d.\ and $\ell_i$ is convex.
        Note that $\nabla \ell_i(0) = v$ and $\nabla \ell_i(\delta) = H \delta + v  = u$.
    \end{enumerate}
    \begin{enumerate}[label=$\clubsuit$, ref=($\clubsuit$)]
        \item\label{item:def_R_j} For indices $j \neq i$, define $\ell_j(\theta) = g_j^\intercal \theta$. 
        Then $\nabla \ell_j(0) = \nabla \ell_j(\delta) = g_j$.
    \end{enumerate}
    Therefore, we get 
    \[
        \F(\delta) = F(\nabla \ell_1(\delta), \ldots, \nabla \ell_n(\delta)) = f_{g_{-i}}(u) \text{ and } \F(0) = F(\nabla \ell_1(0), \ldots, \nabla \ell_n(0)) = f_{g_{-i}}(v).
    \]
    Due to the monotonicity condition~\eqref{eq:operator-convexity}, we have
    \begin{equation}\label{eq:monotonicitycontradiction}
        \langle \F(\delta) - \F(0), \delta - 0 \rangle = \langle f_{g_{-i}}(u) - f_{g_{-i}}(v), \delta \rangle \geq 0.
    \end{equation}
    However, inequality~\eqref{eq:monotonicitycontradiction} contradicts~\eqref{eq:strict_separation}. 
    Thus, for all $u \neq v \in \R^d$, there exists $\lambda(u, v) \geq 0$ such that
    \begin{equation}\label{eq:positivedirectional}
        f_{g_{-i}}(u) - f_{g_{-i}}(v) = \lambda(u,v)(u-v).
    \end{equation}


    \paragraph{Step 2.} We demonstrate that the rigid geometric constraint of preserving input-wise update directions forces the aggregation $F$ to be affine with respect to each individual input vector.
    Evaluating~\eqref{eq:positivedirectional} with $v=0$ yields, for all $u \in \R^d \setminus \{0\}$, $f_{g_{-i}}(u) = \lambda(u, 0)u + f_{g_{-i}}(0)$. 
    We only require to prove that $u \mapsto \lambda(u, 0)$ is a constant. 
    %\Cref{eq:positivedirectional} implies, for all $u \neq v \in \R^d$,
    %\begin{equation}\label{eq:directinaltoaffine-intermediate}
    %    \lambda(u, 0)u - \lambda(v, 0)v = \lambda(u,v)(u-v).
    %\end{equation}
    %Rearranging terms from~\eqref{eq:directinaltoaffine-intermediate}, we have the following linear equation
    %\begin{equation}\label{eq:directinaltoaffine-intermediate2}
    %    (\lambda(u,0) - \lambda(u,v))u + (\lambda(u,v) - \lambda(v,0))v = 0.
    %\end{equation}
    Rearranging the terms in~\eqref{eq:positivedirectional} yields, for all $u \neq v \in \R^d$, the following linear equation
    \begin{equation}\label{eq:directinaltoaffine-intermediate2} 
        (\lambda(u,0) - \lambda(u,v))u + (\lambda(u,v) - \lambda(v,0))v = 0. 
    \end{equation}

    \textit{Case 1: linearly independent inputs.}
    Let any fixed linearly independent vectors $u, v \in \R^d$. 
    By linear independence, the coefficients in~\eqref{eq:directinaltoaffine-intermediate2} must be $0$, i.e., $\lambda(u,0) = \lambda(u,v) = \lambda(v,0)$.

    \textit{Case 2: linearly dependent inputs.}
    If $u = 0 \neq v \in \R^d$,~\eqref{eq:directinaltoaffine-intermediate2} implies $\lambda(v,0) = \lambda(0,v)$.
    Now, let $u \neq v \in \R^d \setminus \{0\}$ be non-zero collinear vectors. 
    Since $d \ge 2$, there exists a vector $z \in \R^d \setminus \{0\}$ linearly independent of $u$, and hence linearly independent of $v$.
    Applying the result from Case 1, $\lambda(u,0) = \lambda(z,0)$ and $\lambda(v,0) = \lambda(z,0)$.
    By transitivity, $\lambda(u,0) = \lambda(v,0)$.

    Consequently, for any pair of vectors $u \neq v \in \R^d $, $\lambda(u,0) = \lambda(v,0)$. 
    Hence $\lambda(\cdot, 0)$ is a constant function.
    Let $\alpha \in \R$ denote this constant, we have proven that 
    \begin{equation}\label{eq:indexaffine-intermediate}
        f_{g_{-i}}(u) = \alpha u + f_{g_{-i}}(0), \text{ for all } u\in\R^d.
    \end{equation}


    \paragraph{Step 3.} We synthesize these input-wise affineness results, proving that the aggregation $F$ must be a globally affine map over the joint domain.
    We derived~\eqref{eq:indexaffine-intermediate} for an arbitrary fixed index $i\in[n]$ and configuration $g_{-i} \in (\R^d)^{n-1}$.
    Therefore, $F$ is input-wise affine, that is, for all $i \in [n]$, all $g_{-i} \in (\R^d)^{n-1}$, we have for all $g_i \in \R^d$,
    \begin{equation}\label{eq:indexaffine}
        F(g_1, \ldots, g_n) = \alpha_i (g_{-i}) g_i + b_i(g_{-i}),
        \quad \text{where } b_i(g_{-i}) \vcentcolon= f_{g_{-i}}(0).
    \end{equation}

    We first show that for any pair of distinct indices $i \neq j \in [n]$, $\alpha_i$ is independent of $g_j$. 
    Without loss of generality, consider indices $1$ and $2$ and fix $\bar{g} \vcentcolon= (g_3, \ldots, g_n) \in (\R^d)^{n-2}$. 
    Equating the representations from~\eqref{eq:indexaffine} yields
    \begin{equation}\label{eq:global-affine-intermediate}
        \alpha_1(g_2, \bar{g}) g_1 + b_1(g_2, \bar{g}) = \alpha_2(g_1, \bar{g}) g_2 + b_2(g_1, \bar{g}).
    \end{equation}
    Evaluating~\eqref{eq:global-affine-intermediate} at $g_1=0$ yields $b_1(g_2, \bar{g}) = \alpha_2(0, \bar{g}) g_2 + b_2(0, \bar{g})$, and further evaluating the relation at $g_2=0$ gives $b_1(0, \bar{g}) = b_2(0, \bar{g})$. 
    Similarly, $b_2(g_1, \bar{g}) = \alpha_1(0, \bar{g}) g_1 + b_1(0, \bar{g})$.
    Substituting this back into~\eqref{eq:global-affine-intermediate}, and rearranging terms to isolate $g_1$ and $g_2$, we have for all $(g_1, \ldots, g_n) \in (\R^d)^n$,
    \begin{equation}\label{eq:lemma_separation}
        (\alpha_1(g_2, \bar{g}) - \alpha_1(0, \bar{g})) g_1 = (\alpha_2(g_1, \bar{g}) - \alpha_2(0, \bar{g})) g_2.
    \end{equation}
    Fix any $g_2 \neq 0$.
    %As~\eqref{eq:lemma_separation} must hold for all $g_1 \in \R^d$, the vector $(\alpha_1(g_2, \bar{g}) - \alpha_1(0, \bar{g}))g_1$ must be collinear with $g_2$ for all $g_1 \in \mathbb{R}^d$.
    %For any $(g_1, \ldots, g_n) \in (\R^d)^n$,~\Cref{eq:lemma_separation} implies that, if $\alpha_1(g_2, g\bar{g}) = \alpha_1(0, \bar{g})$, then $\alpha_2(g_1, \bar{g}) = \alpha_2(0, \bar{g})$, and vice versa.
    %Now assume $\alpha_1(g_2, \bar{g}) \neq \alpha_1(0, \bar{g})$, which forces $\alpha_2(g_1, \bar{g}) \neq \alpha_2(0, \bar{g})$.
    %\Cref{eq:lemma_separation} implies that for any fixed $g_2 \neq 0$, the vector $(\alpha_1(g_2, \bar{g}) - \alpha_1(0, \bar{g}))g_1$ must be collinear with $g_2$ for all $g_1 \in \mathbb{R}^d$.
    Since $\alpha_1(g_2, \bar{g}) - \alpha_1(0, \bar{g})$ is independent of $g_1$, and $d \geq 2$, we can select a vector $g_1'$ linearly independent of $g_2$ in~\eqref{eq:lemma_separation}, so that $(\alpha_1(g_2, \bar{g}) - \alpha_1(0, \bar{g})) g'_1$ is collinear to $(\alpha_2(g'_1, \bar{g}) - \alpha_2(0, \bar{g})) g_2$. 
    The only vector collinear with $(\alpha_2(g'_1, \bar{g}) - \alpha_2(0, \bar{g})) g_2$ in the direction of $g_1'$ is $0$.
    Thus $\alpha_1(g_2, \bar{g}) = \alpha_1(0, \bar{g})$ for all $g_2\in\R^d$, i.e., $\alpha_1$ is independent of $g_2$. 

    By symmetry, for all $(g_1, \ldots, g_n) \in (\R^d)^n$ and each $i \in [n]$, $\alpha_i$ is independent of $g_j$ for all $j \neq i$, and is therefore constant. 
    %(we denote this constant $\alpha_i$, slightly overloading notation).

    Now, define the following residual function, for all $(g_1, \ldots, g_n) \in (\R^d)^n$,
    \begin{equation}\label{eq:residual}
        \textstyle R(g_1, \ldots, g_n) = F(g_1, \ldots, g_n) - \sum_{i=1}^n \alpha_i g_i. 
    \end{equation}
    We prove $R$ is a constant vector.
    Fix $j\in[n]$. Substituting~\eqref{eq:indexaffine} into~\eqref{eq:residual} yields
    \begin{equation}\label{eq:globalaffine-intermediate}
        \textstyle R(g_1, \ldots, g_n) = \alpha_j g_j + b_j(g_{-j}) - \sum_{i=1}^n \alpha_i g_i = b_j(g_{-j}) - \sum_{i \neq j} \alpha_i g_i.
    \end{equation}
    The right-hand side of~\eqref{eq:globalaffine-intermediate} depends only on $g_{-j}$, thus $R$ is independent of $g_j$.
    By symmetry, this holds for every $j \in [n]$, i.e., the residual $R$ is independent of its inputs and is therefore a constant vector.
    Denote this constant by $C \vcentcolon= R(0, \ldots, 0) = F(0, \ldots, 0)$. 
    We have thus shown that
    \begin{equation*}
        \textstyle F(g_1, \ldots, g_n) = \sum_{k=1}^n \alpha_k g_k + C, \text{ for all } (g_1, \ldots, g_n) \in (\R^d)^n.
    \end{equation*}

    Suppose in addition that $F$ is idempotent and, in fact, it suffices that there exist $u \neq v \in \R^d$ such that $F(u, \ldots, u) = u$ and $F(v, \ldots, v) = v$. Then $u-v = (\sum_{i=1}^n\alpha_i) (u-v)$, i.e.\ $\sum_{i=1}^n \alpha_i = 1$. 
    If, moreover, $F$ is permutation-invariant, all coefficients are equal: $\alpha = (\frac{1}{n}, \ldots, \frac{1}{n})^\intercal$.
\end{proof}

\end{document}
